% Part: incompleteness
% Chapter: arithmetization-syntax
% Section: proofs-in-lk

\documentclass[../../../include/open-logic-section]{subfiles}

\begin{document}

\olfileid[pt]{inc}{art}{plk}
\olsection{\usetoken{P}{derivation} em $\Log{LK}$}

\begin{explain}
Para aritmetizar !!{derivation}s, precisamos representar
!!{derivation}s como
números. Como !!{derivation}s são árvores de sequentes nas quais cada
inferência também possui um rótulo, uma representação recursiva é a
abordagem mais evidente: representamos !!a{derivation} como uma tupla
cujos componentes são o sequente final, o rótulo e as representações
das sub-!!{derivation}s que levam às premissas da última inferência.
\end{explain}

\begin{defn}
Se $\Gamma$ é uma sequência finita de !!{sentence}s, $\Gamma =
\tuple{!A_1, \dots, !A_n}$, então $\Gn{\Gamma} = \tuple{\Gn{!A_1},
  \dots, \Gn{!A_n}}$.

Se $\Gamma \Sequent \Delta$ é um sequente, então um número de G\"odel
de $\Gamma \Sequent \Delta$ é
\[
\Gn{\Gamma \Sequent \Delta} = \tuple{\Gn{\Gamma}, \Gn{\Delta}}
\]

Se $\pi$ é !!a{derivation} em $\Log{LK}$, então $\Gn{\pi}$ é definido
da seguinte maneira:
\begin{enumerate}
\item Se $\pi$ consiste apenas no sequente inicial
  $\Gamma \Sequent \Delta$, então $\Gn{\pi}$ é
  \[
    \tuple{0, \Gn{\Gamma \Sequent \Delta}}.
  \]
\item Se $\pi$ termina em uma inferência com uma ou duas premissas, tem
  $\Gamma \Sequent \Delta$ como conclusão, e $\pi_1$ e $\pi_2$ são as
  subprovas imediatas que terminam nas premissas da última inferência,
  então $\Gn{\pi}$ é
  \begin{align*}
    & \tuple{1, \Gn{\pi_1}, \Gn{\Gamma \Sequent \Delta}, k} \text{ ou}\\ 
    & \tuple{2, \Gn{\pi_1}, \Gn{\pi_2}, \Gn{\Gamma \Sequent \Delta}, k}, 
  \end{align*}
  respectivamente, em que $k$ é dado pela tabela a seguir, conforme a
  regra usada na última inferência:

  \begin{tabular}{lccccccc}
    \text{Regra:} & \LeftR{\Weakening} & \RightR{\Weakening} &
    \LeftR{\Contraction} & \RightR{\Contraction} &
    \LeftR{\Exchange} & \RightR{\Exchange} \\
    $k$: & 1 & 2 & 3 & 4 & 5 & 6 \\[2ex]
    \text{Regra:} & \LeftR{\lnot} & \RightR{\lnot} &
    \LeftR{\land} & \RightR{\land} & 
    \LeftR{\lor} & \RightR{\lor} \\
    $k$: & 7 & 8 & 9 & 10 & 11 & 12 \\[2ex]
    \text{Regra:} & \LeftR{\lif} & \RightR{\lif} &
    \LeftR{\lforall} & \RightR{\lforall} &
    \LeftR{\lexists} & \RightR{\lexists} \\
    $k$: & 13 & 14 & 15 & 16 & 17 & 18 \\[2ex]
    \text{Regra:} & \Cut & = \\
    $k$: & 19 & 20
  \end{tabular}
\end{enumerate}
\end{defn}

\begin{ex}
  Considere a !!{derivation} muito simples
  \begin{prooftree}
    \Axiom$!A \fCenter !A$
    \RightLabel{\LeftR{\land}}
    \UnaryInf$!A \land !B \fCenter !A$
    \RightLabel{\RightR{\lif}}
    \UnaryInf$\fCenter (!A \land !B) \lif !A$
  \end{prooftree}
  O número de G\"odel do sequente inicial seria $p_0 = \tuple{0,
    \Gn{!A \Sequent !A}}$. O número de G\"odel da !!{derivation} que
    termina na conclusão de~$\LeftR{\land}$ seria $p_1 =
    \tuple{1, p_0, \Gn{!A \land !B \Sequent !A}, 9}$ ($1$ porque
    $\LeftR{\land}$ tem uma premissa, o número de G\"odel da
    conclusão~$!A \land !B \Sequent !A$, e $9$ é o número que codifica
    $\LeftR{\land}$). O número de G\"odel da !!{derivation} inteira é,
    então, $\tuple{1, p_1, \Gn{\Sequent (!A \land !B) \lif !A)}, 14}$,
    O.s.\@,
  \[
  \tuple{1, \tuple{1, \tuple{0, \Gn{!A \Sequent !A)}}, \Gn{!A \land !B \Sequent !A}, 9},
    \Gn{\Sequent (!A \land !B) \lif !A}, 14}.
  \]
\end{ex}

\begin{explain}
Escolhida uma representação das !!{derivation}s, precisamos também
mostrar que podemos manipular tais !!{derivation}s recursivamente de modo primitivo e
expressar assim suas propriedades e relações essenciais. Algumas
operações são simples: p.ex., dado um número de G\"odel~$p$ de
!!a{derivation}, $\fn{EndSeq}(p) = (p)_{(p)_0+1}$ nos dá o número de
G\"odel de seu sequente final, e $\fn{LastRule}(p) = (p)_{(p)_0+2}$, o
código de sua última regra. A propriedade $\fn{Sequent}(s)$ definida por
\[
  \len{s} = 2 \land \bforall{i<\len{(s)_0} + \len{(s)_1}}{\fn{Sent}(((s)_0 \concat (s)_1)_i)}
\]
vale de $s$ sse $s$ é o número de G\"odel de um sequente constituído de
!!{sentence}s. Algumas propriedades são muito mais difíceis.
Esboçaremos ao menos como tratá-las. O objetivo é mostrar que a relação
``$\pi$ é !!a{derivation} de~$!A$ a partir de~$\Gamma$'' é uma relação
recursiva primitiva dos números de G\"odel de $\pi$ e~$!A$.
\end{explain}

\begin{prop}\ollabel{prop:followsby}
  A propriedade $\fn{Correct}(p)$, que vale sse a última inferência na
  !!{derivation}~$\pi$ com número de G\"odel~$p$ é correta, é recursiva
  primitiva.
\end{prop}

\begin{proof}
  $\Gamma \Sequent \Delta$ é um sequente inicial se existe
  !!a{sentence}~$!A$ tal que $\Gamma \Sequent \Delta$ é
  $!A \Sequent !A$, ou existe um termo~$t$ tal que
  $\Gamma \Sequent \Delta$ é $\emptyset \Sequent \eq[t][t]$. Em termos
  de números de G\"odel, $\fn{InitSeq}(s)$ vale sse
  \begin{align*}
  \bexists{x < s}{} (\fn{Sent}(x) & \land
  s = \tuple{\tuple{x},\tuple{x}})
  \lor {}\\
  \bexists{t<s}{} (\fn{Term}(t) & \land
  s = \tuple{0, \tuple{\Gn{{\eq}(} \concat t \concat \Gn{,} \concat t \concat \Gn{)}}}).
  \end{align*}

  Também precisamos mostrar que, para cada regra de inferência~$R$, a
  relação $\fn{FollowsBy}_R(p)$ é recursiva primitiva, em que
  $\fn{FollowsBy}_R(p)$ vale sse $p$ é o número de G\"odel de
  !!a{derivation}~$\pi$ e o sequente final de~$\pi$ decorre, por uma
  aplicação correta de~$R$, das sub-!!{derivation}s imediatas de~$\pi$.

  Um caso simples é o da regra \RightR{\land}. Se $\pi$ termina em uma
  inferência $\RightR{\land}$ correta, ela tem esta forma:
  \begin{prooftree}
    \AxiomC{}
    \RightLabel{$\pi_1$}
    \Deduce$\Gamma \fCenter \Delta, !A$
    
    \AxiomC{}
    \RightLabel{$\pi_2$}
    \Deduce$\Gamma \fCenter \Delta, !B$

    \RightLabel{\RightR\land}
    \BinaryInf$\Gamma \fCenter \Delta, !A \land !B$
  \end{prooftree}
  Assim, a última inferência na !!{derivation} $\pi$ é uma aplicação
  correta de $\RightR{\land}$ sse existem sequências de !!{sentence}s
  $\Gamma$ e $\Delta$, bem como duas !!{sentence}s $!A$ e~$!B$, tais
  que o sequente final de $\pi_1$ é $\Gamma \Sequent \Delta, !A$, o
  sequente final de $\pi_2$ é $\Gamma \Sequent \Delta, !B$ e o
  sequente final de $\pi$ é $\Gamma \Sequent \Delta, !A \land !B$.
  Basta traduzir isso em números de G\"odel. Se
  $s = \Gn{\Gamma \Sequent \Delta}$, então $(s)_0 = \Gn{\Gamma}$ e
  $(s)_1 = \Gn{\Delta}$. Portanto,
  $\fn{FollowsBy}_{\RightR{\land}}(p)$ vale sse
\begin{align*}
& \bexists{g < p}{\bexists{d < p}{\bexists{a < p}{\bexists{b < p}{\quad}}}} \\
& \qquad \fn{EndSequent}(p) = 
  \tuple{g, d \concat 
    \tuple{\Gn{(} \concat a \concat \Gn{\land} \concat b \concat \Gn{)}}} \land {} \\
& \qquad \fn{EndSequent}((p)_1) = \tuple{g, d \concat \tuple{a}} \land {}\\
& \qquad \fn{EndSequent}((p)_2) = \tuple{g, d \concat \tuple{b}} \land {}\\
& \qquad (p)_0 = 2 \land \fn{LastRule}(p) = 10.
\end{align*}
As linhas individuais expressam, respectivamente: ``existe uma
sequência~($\Gamma$) com número de G\"odel~$g$, existe uma
sequência~($\Delta$) com número de G\"odel~$d$, !!a{formula}~($!A$) com
número de G\"odel~$a$ e !!a{formula}~($!B$) com número de G\"odel~$b$'',
tais que ``o sequente final de $\pi$ é $\Gamma \Sequent \Delta, !A
\land !B$'', ``o sequente final de $\pi_1$ é
$\Gamma \Sequent \Delta, !A$'', ``o sequente final de $\pi_2$ é
$\Gamma \Sequent \Delta, !B$'' e ``$\pi$ tem duas subderivações
imediatas e a última regra de inferência é $\RightR\land$ (com
número~$10$)''.

A última inferência em~$\pi$ é uma aplicação correta de
$\RightR\lexists$ sse existem sequências $\Gamma$ e $\Delta$,
!!a{formula}~$!A$, uma variável~$x$ e um termo~$t$, tais que o sequente
final de $\pi$ é $\Gamma \Sequent \Delta, \lexists[x][!A]$ e o
sequente final de $\pi_1$ é
$\Gamma \Sequent \Delta, \Subst{!A}{t}{x}$. Assim, em termos de números
de G\"odel, temos $\fn{FollowsBy}_{\RightR\lexists}(p)$ sse
\begin{align*}
  & \bexists{g<p}{\bexists{d<p}{\bexists{a<p}{\bexists{x<p}{\bexists{t<p}{\quad}}}}}\\
  & \qquad \fn{EndSequent}(p) = 
  \tuple{
    g,
    d \concat \tuple{\Gn{\lexists} \concat x \concat a}
  } \land {}\\
  & \qquad \fn{EndSequent}((p)_1) = 
  \tuple{
    g,
    d \concat \tuple{\fn{Subst}(a, t, x)}
  } \land {}\\
  & \qquad (p)_0 = 1 \land \fn{LastRule}(p) = 18.
\end{align*}
Definimos então $\fn{Correct}(p)$ como
\begin{multline*}
  \fn{Sequent}(\fn{EndSequent}(p)) \land {}\\ 
  [(\fn{LastRule}(p) = 1 \land 
  \fn{FollowsBy}_{\LeftR\Weakening}(p)) \lor \dots \lor {}\\
  (\fn{LastRule}(p) = 20 \land \fn{FollowsBy}_{\eq}(p)) \lor {}\\
  (p)_0 = 0 \land \fn{InitialSeq}(\fn{EndSequent}(p))]
\end{multline*}
A primeira linha garante que o sequente final de~$d$ é de fato um
sequente constituído de !!{sentence}s. A última linha cobre o caso em
que $p$ é apenas um sequente inicial.
\end{proof}

\begin{prob}
  Defina as propriedades a seguir como em
  \olref[inc][art][plk]{prop:followsby}:
  \begin{enumerate}
  \item $\fn{FollowsBy}_{\Cut}(p)$,
  \item $\fn{FollowsBy}_{\LeftR\lif}(p)$,
  \item $\fn{FollowsBy}_{\eq}(p)$,
  \item $\fn{FollowsBy}_{\RightR{\lforall}}(p)$.
  \end{enumerate}
  Para a última, você também terá de mostrar que é possível testar
  recursivamente de modo primitivo se a última inferência da
  !!{derivation} com número de G\"odel~$p$ satisfaz a condição da
  variável própria, O.s.\@, se a variável própria~$a$ de
  $\RightR{\lforall}$ não ocorre no sequente final.
  \end{prob}
  

\begin{prop}
  \ollabel{prop:deriv}
  A relação $\fn{Deriv}(p)$, que vale se $p$ é o número de G\"odel de
  !!a{derivation} correta~$\pi$, é recursiva primitiva.
\end{prop}

\begin{proof}
  !!^a{derivation}~$\pi$ é correta se cada uma de suas inferências é
  uma aplicação correta de uma regra, O.s.\@, se cada uma de suas
  sub-!!{derivation}s termina em uma inferência correta. Assim,
  $\fn{Deriv}(d)$ sse
  \[
  \bforall{i<\len{\fn{SubtreeSeq}(p)}}{\fn{Correct}((\fn{SubtreeSeq}(p))_i}.
  \]
\end{proof}


\begin{prop}
Suponha que $\Gamma$ seja um conjunto recursivo primitivo de
!!{sentence}s. Então a relação $\Prf[\Gamma](x, y)$ que expressa ``$x$
é o código de !!a{derivation}~$\pi$ de $\Gamma_0 \Sequent !A$ para
algum $\Gamma_0 \subseteq \Gamma$ finito e $y$ é o número de G\"odel
de~$!A$'' é recursiva primitiva.
\end{prop}

\begin{proof}
Suponha que ``$y \in \Gamma$'' seja dada pelo predicado recursivo
primitivo~$R_\Gamma(y)$. Precisamos mostrar que $\Prf[\Gamma](x, y)$,
que vale sse $y$ é o número de G\"odel de uma sentença~$!A$ e $x$~é o
código de uma !!{derivation} em $\Log{LK}$ com sequente final
$\Gamma_0 \Sequent !A$, é recursiva primitiva.

Pela proposição anterior, a propriedade $\fn{Deriv}(x)$, que vale sse
$x$ é o código de !!a{derivation} correta~$\pi$ em $\Log{LK}$, é
recursiva primitiva. Se $x$ é tal código, então $\fn{EndSequent}(x)$ é
o código do sequente final de~$\pi$; portanto,
$(\fn{EndSequent}(x))_0$ é o código do lado esquerdo do sequente final,
e $(\fn{EndSequent}(x))_1$, o do lado direito. Assim, podemos expressar
``o lado direito do sequente final de~$\pi$ é~$!A$'' como
$\len{(\fn{EndSequent}(x))_1} = 1 \land ((\fn{EndSequent}(x))_1)_0 =
x$. O lado esquerdo do sequente final de $\pi$ é, naturalmente,
automaticamente finito; só precisamos expressar que toda sentença nele
pertence a~$\Gamma$. Assim, podemos definir $\Prf[\Gamma](x, y)$ por
\begin{align*}
\Prf[\Gamma](x, y) \defiff {}&
\fn{Deriv}(x) \land {} \\
& \bforall{i <
  \len{(\fn{EndSequent}(x))_0}}{R_\Gamma(((\fn{EndSequent}(x))_0)_i)} \land {}\\
& \len{(\fn{EndSequent}(x))_1} = 1 \land ((\fn{EndSequent}(x))_1)_0 = y.
\end{align*}
\end{proof}

\end{document}
